\section{実験計画法2；一要因モデル(Between)}
\label{lesson12}

\subsection{授業内容}
\subsubsection{科目の中でのこのコマの位置づけ}
\begin{tabular}[t]{ccccccc}
    記述統計[4] & 確率[2] & モデル[6] & 推測・推定[1] & 意思決定[0] & 実践[2]
\end{tabular}
\subsubsection{概要}
実験計画に基づいて，データをどのように扱っていくかを具体的に確認する。まずは一要因一要因3水準，Betweenデザインの数値例を用いる。ここで改めて，標本平均と母平均の違いや平均因果効果の推論をしようとしていることを確認し，統計的に(すなわち母集団において)差があるかないかをどのように考えるかを理解する。回帰分析の時と同じように線形モデルで考え，数式の表現を行い，回帰分析と同じ形式Formであることを確認する。
\subsubsection{コマ主題細目}
\begin{description}
\item[標本平均と母平均]データを元に群の平均を算出し，その大小を比較することは当然可能である。ここに差があることは疑うべくもないが，本来こうした実験計画を行う上で見たかったのは，標本における効果ではなく母集団における効果であり，標本における違いが母集団においても同様に確認できるのか，あるいはたまたまこの標本がこのような違いを生み出しただけなのかは，慎重に議論する必要がある。改めて標本と母集団の関係を考え，仮説は母集団の母数に対してなされることを確認する。
\begin{flushright}
    $\to$母集団と標本について，改めて\citet{Yamada2004}のP.68--73を参照しておく。
\end{flushright}

\item[誤差と効果]RCTデザインをとることで，もし，見ようとする効果がないとするならば，全ての結果変数が全体平均に一致すると考えられる。群に分割することの効果があるなら，それは群内のどの要素に対しても生じているはずであり，群の平均と全体平均の差で算出できるに違いない。こうした理論的な計算に沿って，なお実測値と異なる値が得られるとするならば，測定や実験状況，個人差に伴って生じる誤差になるはずである。このようにして効果と誤差を算出できる。具体的に考えやすくするため，\citet{Yamada2004}のP.162の表7.1.1にある数値例を用いて議論を進める。
\begin{flushright}
    $\to$　\citet{大村201301}のP.64--72が詳しい。\citet{Yamada2004}のP.162--165の数値例と計算式も参照。
\end{flushright}

\item[回帰分析と実験計画]効果と誤差をそれぞれ未知母数$\delta_j,\epsilon_{ij}$とおき，実測値${y_{ij}}$がどのように表現できるかを考える。この時，効果の有無を0と1で表した説明変数$X$を用いることになる。このように数式で表現することで，その形が回帰分析と同じであること，説明変数のデータ種が違うだけであることを再確認できる。
\begin{flushright}
    $\to$ \citet{kosugi2018}のP.136--138
\end{flushright}

\item[デザイン行列の制約]誤差の算出の時とは異なり，回帰分析的表現では切片を第一水準の平均値におき，そこからの相対比較で表現していた。これを全体平均からの相対比較に書き換えることで，両者の数学的表現を同一にすることが可能である。ただしこの場合は，推定できる効果の未知母数が「増える」のではなく，数学的性質から「制約をかける」ことで，実質的に推定すべき母数の数は同じであることが確認できる。


\end{description}
\subsubsection{キーワード}
\begin{itemize}
\item 標本統計量と母数
\item 誤差と効果
\item 回帰分析と実験計画
\item 一般線形モデル
\item 制約
\end{itemize}
\subsection{授業情報}
\paragraph{コマの展開方法}
講義
\paragraph{標準シラバスにおける位置づけ}
科目番号4；心理学研究法；2.データを用いた実証的な思考方法；C データの統計的記述\\
科目番号5；心理学統計法；1.心理学で用いられる統計手法；G 推測統計:代表値と散布度をめぐって(1)
\subsubsection{予習・復習課題}
\paragraph{予習}
統計環境Rを使った線形モデルを実際に行い，切片と傾きが算出されることを確認し，また説明変数が離散的であった場合の切片や傾きがどういう意味を持っているのかを改めて考えておくと理解が進む。
\paragraph{復習}
全体平均，群平均を計算に用いて効果と誤差を算出する過程を，もう一度確認しておく。この時，自分で収集した適当なデータがあると理解が進む。身の回りにあるデータ(ex.お菓子や食べ物の重さなど)や，公開されているデータ(ex.スポーツに関するデータなどは公開されていることが多い)を使って計算するとよい。計算にはExcel/Numbers/Calcsなどの表計算ソフトを用いても良いし，統計環境Rで行列のデータを対象に分析してももちろん良い。


